\subsection{Approximating the time integral form}
Interestingly, the time integral expression of eq. \ref{eq:AGP_time_integral} before taking the limit can be approximated by truncating the infinite integration interval to a finite interval $[-a, a]$:
\begin{align}\label{eq:AGP_time_integral_approx}
I(\omega; \omega_0, a) &:= \frac12 \int_{-a}^a \d\tau \, \e^{-\omega_0|\tau|} \sgn(\tau) \, \e^{-\i \kern0.04em \omega \tau} \nonumber \\
&= \frac{\omega}{\i(\omega^2 + \omega_0^2)} - \e^{-a \+ \omega_0} \frac{\cos(a\omega) \+ \omega + \sin(a\omega) \+ \omega_0}{\i(\omega^2 + \omega_0^2)}.
\end{align}
Note that this expression converges only when we take the limit $a\to\infty$ \emph{before} the limit $\omega_0\to0$. This means that we will always have to work with a nonzero \emph{cutoff gap} $\omega_0$. However, for any nonzero $\omega_0$, the last term in eq. \ref{eq:AGP_time_integral_approx} decays exponentially with time. Taking a series expansion of the second-to-last term in eq. \ref{eq:AGP_time_integral_approx}, we have
\begin{align}
I(\omega; \omega_0, a) = \frac1{\i\omega} + O\Big( \frac{\i\omega_0^2}{\omega^3} \Big) + O(\e^{-\omega_0 a}).
\end{align}
Actually, we can slightly improve the approximation by choosing a Gaussian regulator instead, i.e.
\begin{align} \label{eq:gaussian_regulator}
    I(\omega; \omega_0, a) &:= \frac12 \int_{-a}^a \d\tau \, \e^{-\omega_0^2\tau^2} \sgn(\tau) \, \e^{-\i \kern0.04em \omega \tau} \nonumber \\[1mm]
    &= -\frac{\i \sqrt{\pi}}{2\omega_0}\e^{-\frac{\omega ^2}{4\omega_0^2}} \erfi \Big( \frac{\omega}{2 \omega_0} \Big) + O(\e^{-\omega_0^2 a^2}) \nonumber \\[1mm]
    &= \frac1{\i\omega} + O \Big( \frac{\i\omega_0^2}{\omega^3} \Big) + O(\e^{-\omega_0^2 a^2}). \\[2mm]
    & {\color{red} \text{ [DOUBLE CHECK STUFF IN O NOTATION]}} \nonumber
\end{align}
A digitised form of the CD-assisted adiabatic evolution can be implemented as a Trotter product
\begin{align}\label{eq:CD_Trotter}
\mt U_{\rm CD} &:= \mathcal T \exp \bigg[ \! -\!\i \int_{t_{\rm \i}}^{t_{\rm f}} \d t' \, \mt H_{\rm CD}(\lambda(\d t')) \bigg] \nonumber \\
&\approx \prod_k \exp[-\i \+ \mt H(\lambda(k\+\delta t)) \+ \delta t] \exp[-\i \+ \dot\lambda(k\+\delta t) \+ \mt A(\lambda(k\+\delta t)) \+ \delta t].
\end{align}
In fact, since time is meaningless in the limit of transitionless driving, we can omit $\e^{-\i\+\mt H\+\delta t}$ altogether from $\mt U_{\rm CD}$. 
Then, the approximate integral form of $\mt A$ may itself be approximated as a Riemann sum
\begin{align}
    \frac12 \int_{-a}^a \d\tau \, \e^{-\omega_0^2\tau^2} \sgn(\tau) \, \e^{-\i \mt H \tau} \partial_\lambda \mt \+ \e^{\i \mt H \tau} \approx \frac12 \sum_\kappa \delta\tau \, \e^{-\omega_0^2\tau^2} \sgn(\tau) \, \e^{-\i \mt H \tau} \partial_\lambda \mt H \+ \e^{\i \mt H \tau}
\end{align}
so that we may apply another Trotter expansion to obtain
\begin{align}
\exp[-\i \+ \dot\lambda \+ \mt A \+ \delta t] &\approx \exp \bigg[ -\!\i \+ \dot\lambda \+ \bigg( \frac12 \int_{-a}^a \d\tau \, \e^{-\omega_0^2\tau^2} \sgn(\tau) \, \e^{-\i \mt H \tau} \partial_\lambda \mt H \+ \e^{\i \mt H \tau} \bigg) \delta t \bigg] \nonumber \\
&\approx \prod_\kappa \exp \bigg[ \e^{-\i\mt H \kappa \kern0.04em \delta\tau} \bigg( \! -\!\frac\i2 \+ \dot\lambda \+ \delta t \+ \delta\tau \+ \e^{-\omega_0^2(\kappa\delta\tau)^2} \sgn(\kappa\delta\tau) \+ \partial_\lambda\mt H \bigg) \e^{\i\mt H \kappa\kern0.04em \delta\tau} \bigg] \nonumber \\
&= \prod_\kappa \exp[-\i \+ \mt H \kern0.04em \kappa \+ \delta\tau] \exp \bigg[ \! -\!\frac\i2 \+ \dot\lambda \+ \delta t \+ \delta\tau \+ \e^{-\omega_0^2(\kappa\delta\tau)^2} \sgn(\kappa\delta\tau) \+ \partial_\lambda\mt H \bigg] \exp[\i \+ \mt H \kern0.04em \kappa \+ \delta\tau]
\end{align}
where one should read $\mt H = \mt H(\lambda)$, $\partial_\lambda\mt H = \partial_\lambda\mt H(\lambda)$, $\mt A = \mt A(\lambda)$, $\lambda = \lambda(k\+\delta t)$, $\dot\lambda = \dot\lambda(k\+\delta t)$. Each factor in this expansion can be implemented via standard hamiltonian simulation techniques.
\subsection{Error and infidelity upper bounds}
Now the question is what fidelity can be achieved with this approximation to the AGP, in terms of the coefficient error $|I(\omega; \omega_0, a) - 1/\i\omega|$. 
With this relaxation, we can use the triangle inequality to bound the error as the sum of multiple contributions. We split the error in the integrand in three parts: (i) the finite-time truncation error, (ii) the regularisation error and (iii) the Riemann-Trotter error. \par
If we insert the approximated time-integral form with a Gaussian regulator (see eq. \ref{eq:gaussian_regulator}) into the integrand of eq. \ref{eq:infidelity_bound_norm}, the error density from finite-time truncation can be straightforwardly estimated:
\begin{align}
    \text{truncation error density} &\leq \frac12 \bigg\| \int_a^\infty \d\tau \, \e^{-\omega_0^2 \tau^2}  \e^{-\i \mt H \tau} \, \partial_\lambda \mt H \, \e^{\i \mt H \tau} - \int_{-\infty}^{-a} \d\tau \, \e^{-\omega_0^2 \tau^2} \e^{-\i \mt H \tau} \, \partial_\lambda \mt H \, \e^{\i \mt H \tau} \bigg\| \nonumber \\
    &\leq \int_a^\infty \d\tau \, \|\e^{-\omega_0^2 \tau^2} \e^{-\i \mt H \tau} \, \partial_\lambda \mt H \, \e^{\i \mt H \tau}\| \nonumber \\
    &= \int_a^\infty \d\tau \, \e^{-\omega_0^2 \tau^2} \|\partial_\lambda \mt H\| \nonumber \\
    &= \frac{\sqrt\pi}{2\omega_0}\erfc(\omega_0 a) \|\partial_\lambda \mt H\| \nonumber \\
    &\leq \frac1{2 \omega_0^2 a} \e^{-\omega_0^2 a^2} \|\partial_\lambda \mt H\|
\end{align}
(see the work of Karagiannidis et al. \cite{Karagiannidis2007} for the bound in the last line), and thus
\begin{align}
    \text{truncation error} \leq \frac1{2 \omega_0^2 a} \e^{-\omega_0^2 a^2} \int_{\lambdai}^{\lambdaf} \d\lambda \, \|\partial_\lambda \mt H\|.
\end{align}\par
The regularisation error can also be upper bounded with a relatively simple calculation, using eq. \ref{eq:infidelity_vector_norm}. For any instantanous eigenstate $\ket{k(\lambda)}$, we wish to upper bound the vector norm
\begin{align}
    \text{regularisation error density} = \bigg\| \int_{-\infty}^\infty \d\tau \sgn(\tau) \+ (1 - \e^{-\omega_0^2\tau^2}) \+ \e^{-\i\mt H\tau} \partial_\lambda \mt H \+ \e^{\i\mt H\tau} \ket k \bigg\|.
\end{align}
To make this explicit, we first define the truncated regularised AGP $\tilde{\mt A}_{\rm tr}$ with a general regulariser $r_{\omega_0}(\tau)$ as 
\begin{align}
    \tilde{\mt A}_{\rm tr} &= \int_{-a}^a \d\tau \, \sgn(\tau) \, r_{\omega_0}(\tau) \, \e^{-\i\mt H\tau} \partial_\lambda \mt H \+ \e^{\i\mt H\tau} \nonumber \\
    &= \int_{-\infty}^\infty \d\tau \, \sgn(\tau) \, r_{\omega_0}^a(\tau) \, \e^{-\i\mt H\tau} \partial_\lambda \mt H \+ \e^{\i\mt H\tau}
\end{align}
 To use eq. \ref{eq:infidelity_vector_norm}, we make the energy eigenbasis expansion
\begin{align} \label{eq:time_translated_dlambdaH_energy_eigenbasis}
    \e^{-\i\mt H\tau} \partial_\lambda \mt H \+ \e^{\i\mt H\tau}
    %&= \sum_n \partial_\lambda E_n \ket n\bra n - \sum_{n, \, m\neq n} \e^{-\i (E_m - E_n) \tau} (E_m - E_n) \ket m \braket m{\partial_\lambda n} \bra n \nonumber \\
    &= \sum_n \partial_\lambda E_n \ket n\bra n + \sum_{n, \, m\neq n} \e^{-\i (E_m - E_n) \tau} \ket m \bra m \partial_\lambda \mt H \ket n \bra n.
\end{align}
$\tilde{\mt A}_{\rm rt}$ will respect the gauge if we take $r_{\omega_0}(\tau)$ to be a symmetric function of $\tau$, since then
\begin{align}
    \sum_n \int_{-\infty}^\infty \d\tau \sgn(\tau) \+ (1 - r_{\omega_0}^a(\tau)) \+ \partial_\lambda E_n \ket n \bra n = 0.
\end{align}
Furthermore, for any eigenstate $\ket k$,
\begin{align}
    & \sum_{n, \, m\neq n} \int_{-\infty}^\infty \d\tau \sgn(\tau) \+ (1 - r_{\omega_0}^a(\tau)) \+ \e^{-\i (E_m - E_n) \tau} \ket m \bra m \partial_\lambda \mt H \ket n \braket nk \nonumber \\
    =& \sum_{m\neq k} \Big( \frac1{\i \+ \omega_{mk}} - \hat r_{\omega_0}^a(\omega_{mk}) \Big) \bra m \partial_\lambda \mt H \ket k \+ \ket m
\end{align}
where $\hat r_{\omega_0}^a$ is the Fourier transform of $r_{\omega_0}^a$, so that
\begin{align}
    &\int_{\lambda_{\rm \i}}^{\lambdaf} d\lambda \, \|(\mt A(\lambda) - \tilde{\mt A}_{\rm rt}(\lambda)) \ket{n(\lambda)}\| \nonumber \\
    \leq& \int_{\lambda_{\rm \i}}^{\lambdaf} d\lambda \, \sqrt{\Big| \frac1{\i \+ \omega_{mk}} - \hat r_{\omega_0}^a(\omega_{mk}) \Big|^2 |\bra k\partial_\lambda\mt H\ket m|^2}.
\end{align}
where $I(\omega; \omega_0) := I(\omega; \omega_0, \infty)$. The (non-normalised) state $\ket\Psi$ has norm \footnote{The question marked inequalities need a final check though I'm quite certain they're correct, see eq. \ref{eq:gaussian_regulator}.}
\begin{align}
    \|\ket\Psi\| &= \sqrt{ \sum_{m\neq k} \Big| \frac1{\i \+ \omega_{mk}} - I(\omega_{mk}; \omega_0) \Big|^2 \bra k\partial_\lambda\mt H\ket m\bra m\partial_\lambda\mt H\ket k } \nonumber \\
    &\questionover\leq \sqrt{ \sum_{m\neq k} \Big( \frac{4\omega_0^2}{\omega_{mk}^3} \Big)^2 \bra k\partial_\lambda\mt H\ket m\bra m\partial_\lambda\mt H\ket k } \nonumber \\
    &\leq \frac{4\omega_0^2}{\Delta_k^3} \|\partial_\lambda\mt H\|
\end{align}
where $\Delta_k$ is the minimum absolute gap around $\ket k$; therefore
\begin{align}
    \text{regularisation error} \leq \frac{4\omega_0^2}{\Delta_k^3} \int_{\lambdai}^{\lambdaf} \d\lambda \, \|\partial_\lambda\mt H\|.
\end{align}

% To use this formula, one needs to upper bound the AGP error $\|\mt A(\lambda) - \tilde{\mt A}(\lambda)\|$ from the coefficient error. This also requires knowledge about the scaling of the matrix elements $\bra m\partial_\lambda\mt H\ket n$. According to Kolodrubetz et al. \cite{Kolodrubetz2017}, these matrix elements scale as $\exp[-S/2]$ where $S$ is the extensive thermodynamic entropy, by the eigenstate thermalisation hypothesis. (Again though, this is vague and probably nonrigorous in a typical physicsy way.)\par

Next, there's the Riemann error {\color{red} [ELABORATE]}. We use again eqs. \ref{eq:infidelity_vector_norm} and \ref{eq:time_translated_dlambdaH_energy_eigenbasis} to find
\begin{align}
    \text{Riemann err. dens.} &\leq \bigg\| \bigg( \int_{-a}^a \d\tau \sgn(\tau) \, \e^{-\omega_0^2\tau^2} \e^{-\i\mt H\tau} \partial_\lambda \mt H \+ \e^{\i\mt H\tau} - \sum_\kappa \sgn(\tau_\kappa) \, \e^{-\omega_0^2\tau_\kappa^2} \e^{-\i\mt H\tau_\kappa} \partial_\lambda \mt H \+ \e^{\i\mt H\tau_\kappa} \bigg) \ket n \bigg\| \nonumber \\
    &= \bigg\| \sum_{m\neq n} \bra m \partial_\lambda \mt H \ket n \bigg( \int_{-a}^a \d\tau \sgn(\tau) \, \e^{-\omega_0^2\tau^2} \e^{-\i \+ \omega_{mn}\tau} - \sum_\kappa \delta\tau \sgn(\tau_\kappa) \, \e^{-\omega_0^2\tau_\kappa^2} \e^{-\i \+ \omega_{mn}\tau_\kappa} \bigg) \ket m \bigg\| \nonumber \\
    &\questionover\leq \sqrt{ \sum_{m\neq n} |\bra n \partial_\lambda\mt H \ket m|^2 \, |\delta\tau \+ \e^{-\omega_0^2 a^2}(\e^{-\i \+ \omega_{mn}a} + \e^{\i \+ \omega_{mn}a})|^2 } \nonumber \\
    &\leq 2 \+ \delta\tau \+ \e^{-\omega_0^2 a^2} \|\partial_\lambda\mt H\|
\end{align}
where the second inequality follows from the fact that the Riemann error is upper bounded by the contributions at the boundaries: $\big| \! \int_a^b \d x \+ f(x) - \sum_\kappa \delta x \+ f(x_\kappa) \big| \leq \delta x |f(a) - f(b)|$ {\color{red} [double check this for complex integrals and clear up definition of $\tau_\kappa$]}. Thus
\begin{align}
    \text{Riemann error} \questionover\leq \delta\tau \, \e^{-\omega_0^2 a^2} \int_{\lambdai}^{\lambdaf} \d\lambda \, \|\partial_\lambda\mt H\|.
\end{align}

To be determined are the error from the Riemann sum approximation (this should still respect the gauge) the time-independent Trotter error from implementing the exponential of the Riemann sum {\color{red} which may in fact be reduced by using non-Trotter methods} and a time-dependent Trotter error from the implementation of the sweep. 

\subsection{First idea}
{\color{red} I'm quite certain the proposed ansatz here is not complete, and it already scales pretty badly in the number of modes.}\par
A {\color{red}(not so)} general gauge potential ansatz we can take for such systems, in order to construct the counterdiabatic hamiltonian, is the two-body operator
\begin{align}
    \mt X = \i\sum_{ijkl} \chi_{ijkl} \sum_{\alpha\beta} \atataa{\i\alpha}{k\beta}{l\beta}{j\alpha}.
\end{align}
Note that since $\mt H_{\rm CD}$ is real, $\mt X$ must be imaginary. Furthermore is must be hermitian, which requires that
\begin{align}
    \chi_{ijkl} = -\chi_{klij} \in \R.
\end{align}
For this case, the action minimisation condition reads
\begin{align}
    \paf{}{\chi_{ijkl}} \big( ( 2\mt V + [\i\mt X, \mt H(\lambda)] ) [\i\mt X, \mt H(\lambda)] \big) = 0.
\end{align}
Now since $[\mt X, \mt H]$ is a three-body operator, writing down the action $S$ requires taking traces of six-body operators. We are helped here by the fact that only those six-body operators which can be written as a string of number operators have nonzero trace. Writing $(PP')$ as a shorthand for $\ata P{P'}$, we have
\begin{align}
    \tr[(PP')(QQ')(RR')(SS')(TT')(UU')] = \begin{cases}
        \binom{L - 6}{N - 6} & \text{if } \pi(P'Q'R'S'T'U') = PQRSTU \text{ for some permutation $\pi\in S_6$}, \\
        0 & \text{ otherwise}
    \end{cases}
\end{align}
since any two creation operators in $\mt G_\lambda(\mt X)^2$ can be freely exchanged, and similarly for annihilation operators.
As such, if $\mt G_\lambda^2(\mt X)$ is some six-body operator
\begin{align}
    \mt G_\lambda^2(\mt X) = \sum_{\substack{(PQRSTU) \\ (P'Q'R'S'T'U')}} g_{(PQRSTU)(P'Q'R'S'T'U')} (PP')(QQ')(RR')(SS')(TT')(UU')
\end{align}
then the action minimisation equation (eq. \ref{eq:dSdX=0}) becomes
\begin{align}
    \frac\partial{\partial \chi_{ijkl}} \sum_{\pi\in S_6} g_{(PQRSTU)\pi(PQRSTU)} = 0.
\end{align}
Since $g_{(PQRSTU)}$ is quadratic in the $\chi$ coefficients, this is a linear equation that should be soluble (numerically) through conventional methods, in time polynomial in $L$. This seems too good to be true...\par
Another option, of course, is to pick a less general ansatz.\par
Let's see where we can take this...

\section{Open questions}
Wat zegt de fout $S = \tr[G_\lambda^2]$ over hoe accuraat je geoptimaliseerde ansatz is? Relatie met fidelity via fourierafstand en MT-grens?\par
Geeft benaderde CD een monotone verbetering?\par
Simulaties met optimalisatieproblemen (klassieke eindhamiltonianen)
\begin{itemize}
    \item Is het zo dat CD verstrengeling in het midden van het pad vermindert?
    \item Wordt de energiekloof opengemaakt door CD?
\end{itemize}
QAOA
\begin{itemize}
    \item Hoe kunnen we dit gebruiken om QAOA te verbeteren?
\end{itemize}



\subsection{Sampling approach}
The following is a method inspired by the qDRIFT protocol \cite{}. The approach still discretises $\lambda$ space, in that we approximate $\mt U_\eta$ as a product of unitaries
\begin{align}
    \mt U_\eta(\lambdaf, \lambdai) \approx \tilde{\mt U}_\eta(\lambdaf, \lambdai) = \prod_{j=1}^r \tilde{\mt U}_\eta(\lambda_j, \lambda_{j-1})
\end{align}
such that $\lambda_j > \lambda_{j-1}$, $\lambda_0 = \lambdai$ and $\lambda_r = \lambdaf$. Contrary to the deterministic approach, the $\tilde{\mt U}_j$ are now sampled according to an appropriate distribution. The idea is based on a Monte Carlo evaluation of the time integral expression for $\mt A_\eta(\lambda_j)$. In this setting, one would sample
\begin{align}
    \mt B_\eta(\lambda, \tau) = \e^{-\i \mt H(\lambda) \tau} \Big( \frac{\sgn\tau}\eta \partial_\lambda \mt H(\lambda) \Big) \> \e^{\i \mt H(\lambda) \tau}, \quad P(\tau) = \frac\eta2 \e^{-\eta|\tau|}
\end{align}
so that the expected value of $\mt B(\lambda, \tau)$ over $\tau$ equals $\mt A_\eta(\lambda)$. We port this idea to unitary evolution under $\mt A_\eta(\lambda)$; that is to say, we construct the channel
\begin{align}
    \Ut_\eta(\lambda_j, \lambda_{j-1})(\rhom) = \int_{\lambda_{j-1}}^{\lambda_j} \d\lambda \> p(\lambda) \int_{-\infty}^{\infty} \d\tau P(\tau) \, \tilde{\mt U}_\eta(\lambda, \tau) \> \rhom \> \tilde{\mt U}_\eta\t(\lambda, \tau)
\end{align}
for some distribution $p(\lambda)$ to be specified later, where
\begin{align}
    \tilde{\mt U}_\eta(\lambda, \tau) &= \exp[-\i \mt B_\eta(\lambda, \tau) / p(\lambda)] \nonumber \\
    &= \exp[-\i \mt H(\lambda) \tau] \exp \! \bigg[ \!-\!\i \> \frac{\sgn\tau}{p(\lambda)\eta} \> \partial_\lambda \mt H(\lambda)\> \bigg] \exp[\i \mt H(\lambda)\tau].
\end{align}

\subsubsection{Error bound}
The error we now want to bound is
\begin{align}\label{eq:error_sampling_channel}
    \bigg\| \, \U_\eta(\lambdaf, \lambdai)(\rhom_n(\lambdai)) - \prod_{j=1}^r \Ut_\eta(\lambda_j, \lambda_{j-1})(\rhom_n(\lambdai)) \bigg\|_1.
\end{align}
where $\U_\eta(\lambdaf, \lambdai)(\rho) = \mt U_\eta(\lambdaf, \lambdai) \rhom \mt U_\eta\t(\lambdaf, \lambdai)$ and the product sign is meant to indicate composition. In a similar style to eq. \ref{eq:telescoping_proof_12}, one may telescope this error into multiple short evolution errors:
\begin{align}
    \text{(eq. \ref{eq:error_sampling_channel})} \leq \sum_{j=1}^r \|\> \U_\eta(\lambda_j, \lambda_{j-1})(\rhom_n(\lambda_{j-1})) - \Ut_\eta(\lambda_j, \lambda_{j-1})(\rhom_n(\lambda_{j-1}))\|_1.
\end{align}
We will now bound each term following a similar strategy to that of Berry et al. \cite{Berry2020}. First we define scaled versions of $\U_\eta$ and $\Ut_\eta$:
\begin{align}
    \U_{\eta, s}(\lambda_j, \lambda_{j-1})(\rhom) &= \mt U_{\eta, s}(\lambda_j, \lambda_{j-1}) \rhom \mt U_{\eta, s}\t(\lambda_j, \lambda_{j-1}), \nonumber \\
    \mt U_{\eta, s}(\lambda_j, \lambda_{j-1}) &= \mathcal P \exp \bigg[ \! -\!\i \int_{\lambda_{j-1}}^{\lambda_j} \d\lambda \> s \> \mt A_\eta(\lambda) \bigg] ; \\[1mm]
    \Ut_{\eta, s}(\lambda_j, \lambda_{j-1})(\rhom) &= \int_{\lambda_{j-1}}^{\lambda_j} \d\lambda \> p(\lambda)  \int_{-\infty}^{\infty} \d\tau P(\tau) \, \tilde{\mt U}_{\eta, s}(\lambda, \tau) \> \rhom \> \tilde{\mt U}_{\eta, s}\t(\lambda, \tau), \nonumber \\
    \tilde{\mt U}_{\eta, s}(\lambda, \tau) &= \exp[-\i \> s \> \mt B_\eta(\lambda, \tau) / p(\lambda)].
\end{align}
Since $\U_{\eta, 0}(\rhom) = \Ut_{\eta, 0}(\rhom) = \rhom$, we may write (dropping the $\lambda_j, \lambda_{j-1}$ arguments for legibility)
\begin{align}
    \|\> \U_\eta(\rhom_n) - \Ut_\eta(\rhom_n)\|_1 &= \big\|\> \big\{ \U_{\eta, 1}(\rhom_n) - \U_{\eta, 0}(\rhom_n) \big\} - \big\{ \Ut_{\eta, 1}(\rhom_n) - \Ut_{\eta, 0}(\rhom_n) \big\} \big\|_1 \nonumber \\
    &= \bigg\| \int_0^1 \d s \frac\d{\d s} \big( \> \U_{\eta, s}(\rhom_n) - \Ut_{\eta, s}(\rhom_n) \big) \bigg\|.
\end{align}
For the derivatives in the last line, we use lemma 6 from Berry et al. \cite{Berry2020},
\begin{align}
    \frac\d{\d s} \mt U_{\eta, s}(\lambda_j, \lambda_{j-1}) = \int_{\lambda_{j-1}}^{\lambda_j} \d\lambda \, \mt U_{\eta, s}(\lambda_j, \lambda) [-\i \mt A_\eta(\lambda)] \mt U_{\eta, s}(\lambda, \lambda_{j-1}),
\end{align}
and
\begin{align}
    \frac\d{\d s} \Ut_{\eta, s}(&\lambda_j, \lambda_{j-1})(\rhom) = \int_{\lambda_{j-1}}^{\lambda_j} \d\lambda \> p(\lambda) \int_{-\infty}^\infty \d\tau \, \frac\eta2 \e^{-\eta|\tau|} \nonumber \\[1mm]
    &\times \e^{-\i \mt H(\lambda) \tau} \e^{-\i \> s \sgn\tau \> \partial_\lambda \mt H(\lambda) / p(\lambda) \eta} \nonumber \\
    &\times \Big[ \!-\!\<\i \> \frac{\sgn\tau}{p(\lambda) \eta} \> \partial_\lambda \mt H(\lambda), \, \e^{\i \mt H(\lambda) \tau} \rhom \, \e^{-\i \mt H(\lambda) \tau} \Big] \nonumber \\
    &\times \e^{\i \> s \sgn\tau \> \partial_\lambda \mt H(\lambda) / p(\lambda) \eta} \e^{\i \mt H(\lambda) \tau}.
\end{align}
We observe now that the derivatives of $\U_{\eta, s}$ and $\Ut_{\eta, s}$ agree at $s=0$:
\begin{align}
    \frac\d{\d s} \U_{\eta, s}(\lambda_j, \lambda_{j-1})(\rhom) \Big|_{s=0} = \int_{\lambda_{j-1}}^{\lambda_j} \d\lambda \, [-\i \mt A_\eta(\lambda), \rhom] = \frac\d{\d s} \Ut_{\eta, s}(\lambda_j, \lambda_{j-1})(\rhom) \Big|_{s=0} \, ;
\end{align}
hence, we may invoke the fundamental theorem of calculus a second time to express
\begin{align}
    \|\> \U_\eta(\rhom_n) - \Ut_\eta(\rhom_n)\|_1 = \bigg\| \int_0^1 \d s \int_0^s \d v \frac{\d^2}{\d v^2} \big( \> \U_{\eta, v}(\rhom_n) - \Ut_{\eta, v}(\rhom_n) \big) \bigg\|.
\end{align}
We proceed to bound this error following Berry et al. \cite[proof of theorem 5]{Berry2020}:
\begin{align}\label{eq:sampling_error_proof}
    \|\> \U_\eta(\rhom) - \Ut_\eta(\rhom)\|_1 &\leq \int_0^1 \d s \int_0^s \d v \, \bigg\{ \Big\| \frac{\d^2}{\d v^2} \mt U_{\eta, v} \> \rhom \> \mt U_{\eta, v}\t \Big\|_1 + 2 \Big\| \frac\d{\d v} \mt U_{\eta, v} \> \rhom \> \frac\d{\d v} \mt U_{\eta, v}\t \Big\|_1 + \Big\| \mt U_{\eta, v} \> \rhom \> \frac{\d^2}{\d v^2} \mt U_{\eta, v}\t \Big\|_1 \nonumber \\
    &\hspace{6mm} + \int_{\lambda_{j-1}}^{\lambda_j} \d\lambda \, p(\lambda) \int_{-\infty}^\infty \d\tau \> P(\tau) \nonumber \\
    &\hspace{10mm} \times \Big\| \Big[ \!-\!\<\i \> \frac{\sgn\tau}{p(\lambda)\eta} \> \partial_\lambda \mt H(\lambda), \, \Big[ \!-\!\<\i \> \frac{\sgn\tau}{p(\lambda)\eta} \> \partial_\lambda \mt H(\lambda), \, \e^{\i \mt H(\lambda) \tau} \rhom \, \e^{-\i \mt H(\lambda) \tau} \Big] \Big] \Big\|_1 \bigg\} \nonumber \\
    &\leq \int_0^1 \d s \int_0^s \d v \, \bigg\{ 4 \bigg( \int_{\lambda_{j-1}}^{\lambda_j} \d\lambda \, \|\mt A_\eta(\lambda)\|\bigg)^2 + 4 \int_{\lambda_{j-1}}^{\lambda_j} \d\lambda \, p(\lambda) \bigg( \frac{\|\partial_\lambda \mt H(\lambda)\|}{p(\lambda)\eta} \bigg)^2 \bigg\}.
\end{align}
Now from the definition of $\mt A_\eta$, we find
\begin{align}
    \|\mt A_\eta\| \leq \int_{-\infty}^{\infty} \d\tau \> \frac12\e^{-\eta|\tau|} \|\partial_\lambda \mt H\| = \frac1\eta \|\partial_\lambda \mt H\|
\end{align}
and if we conveniently choose
\begin{align}
    p(\lambda) = \|\partial_\lambda \mt H(\lambda)\| \bigg/ \!\! \int_{\lambda_{j-1}}^{\lambda_j} \d\lambda \, \|\partial_\lambda \mt H(\lambda)\|,
\end{align}
then the last line of eq. \ref{eq:sampling_error_proof} reduces to
\begin{align}
    \|\> \U_\eta(\lambda_j, \lambda_{j-1})(\rhom) - \Ut_\eta(\lambda_j, \lambda_{j-1})(\rhom)\|_1 \leq \frac4{\eta^2} \bigg( \int_{\lambda_{j-1}}^{\lambda_j} \d\lambda \, \|\partial_\lambda \mt H(\lambda)\| \bigg)^2.
\end{align}
Lastly, we can directly follow the proof of theorem 7 from Berry et al. \cite{Berry2020} to show that
\begin{align}
    r \geq \bigg\lceil \frac4{\eta^2\epsilon} \bigg( \int_\lambdai^\lambdaf \d\lambda \, \|\partial_\lambda \mt H(\lambda)\| \bigg)^{\!2} \, \bigg\rceil
\end{align}
is sufficient to guarantee that the error from eq. \ref{eq:error_sampling_channel} satisfies
\begin{align}
    \bigg\| \, \U_\eta(\lambdaf, \lambdai)(\rhom_n(\lambdai)) - \prod_{j=1}^r \Ut_\eta(\lambda_j, \lambda_{j-1})(\rhom_n(\lambdai)) \bigg\|_1 \leq \epsilon.
\end{align}

\subsubsection{Complexity considerations}
According to the error analysis above, one should sample and apply unitary
\begin{align}
    \tilde{\mt U}_\eta(\lambda, \tau) = \exp[-\i \mt H(\lambda) \tau] \exp \! \bigg[ \!-\!\i \> \frac{\sgn\tau}{p(\lambda)\eta} \> \partial_\lambda \mt H(\lambda)\> \bigg] \exp[\i \mt H(\lambda)\tau]
\end{align}
at least $\sim \eta^{-2}\epsilon^{-1}$ times to limit the total error to $O(\epsilon)$. The gate complexity of an implementation of this unitary, however, will cost $O(\|\mt H(\lambda)\|\tau + \|\partial_\lambda \mt H(\lambda)\| / p(\lambda) \eta)$. This means that the total gate complexity will scale as $\sim \eta^{-3}$. Since we already established that $\eta \sim \Delta^{3/2}$, this results in a gate complexity scaling of $\sim \Delta^{-9/2}$, which is worse than the desired $\Delta^{-3}$. Also note that the variance of $\tau$ is $2 / \eta^2$ when sampled from the distribution $P(\tau) = \eta \> \e^{-\eta|\tau|} / 2$. Not quite sure what this entails but it will probably show up in the number of shots needed to make an accurate measurement of some operator after the desired state has been prepared.